\chapter{リリース管理 --- ソフトウェアを世に届ける}

\section{この章で学ぶこと}

この章では、GitHubにおけるソフトウェアの\textbf{リリース管理}について学びます。具体的には、以下の内容を扱います。

\begin{itemize}
  \item タグ（Tag）とは何か――特定のコミットに名前を付ける仕組み
  \item 軽量タグと注釈付きタグの違いと使い分け
  \item セマンティックバージョニング（SemVer）――バージョン番号の付け方の標準的な規約
  \item GitHubリリースの作成方法（Web UIとCLIの両方）
  \item リリースノートの自動生成とカスタマイズ
\end{itemize}

ソフトウェア開発は「作って終わり」ではありません。作ったものをユーザーに届け、バージョンごとに何が変わったかを伝え、過去のバージョンに戻れるようにしておくことが、信頼されるソフトウェアの条件です。リリース管理は、その「届ける」プロセスを体系的に行うための仕組みです。

\section{タグ（Tag）とは}

\subsection{コミットに名前を付ける}

第9章で学んだように、Gitのコミットは \cmd{a1b2c3d} のようなハッシュ値で識別されます。しかし、「このハッシュ値がバージョン1.0のリリースだった」ということを覚えておくのは困難です。そこで使うのが\textbf{タグ（Tag）}です。

タグは、特定のコミットに対して人間が読みやすい名前（多くの場合バージョン番号）を付ける機能です。コミット履歴の中に「しおり」を挟むようなものだと考えてください。

\begin{lstlisting}[language={}]
A → B → C → D → E → F (HEAD)
         ↑       ↑
       v1.0.0   v2.0.0   ← タグ
\end{lstlisting}

上の図では、コミットCに \cmd{v1.0.0}、コミットEに \cmd{v2.0.0} というタグが付けられています。タグは後から追加することもでき、コミット履歴を変更するものではありません。あくまで「このコミットがこのバージョンに相当する」という\textbf{目印}です。

\subsection{軽量タグと注釈付きタグ}

Gitには2種類のタグがあります。

\textbf{軽量タグ（Lightweight Tag）}は、特定のコミットを指すポインタだけを作成する最もシンプルなタグです。ブランチのように「名前がついた参照」を作るだけで、タグ自体に付加情報は持ちません。

\begin{lstlisting}[language=bash]
# 軽量タグの作成
git tag v1.0.0
\end{lstlisting}

\textbf{注釈付きタグ（Annotated Tag）}は、タグに対してメッセージ（注釈）、作成者の情報、作成日時を記録できるタグです。Gitの内部ではタグオブジェクトとして保存されるため、「誰が、いつ、なぜこのバージョンを付けたか」という情報が残ります。

\begin{lstlisting}[language=bash]
# 注釈付きタグの作成（-a フラグと -m でメッセージを指定）
git tag -a v1.0.0 -m "Version 1.0.0 - Initial stable release"
\end{lstlisting}

リリースに使うタグとしては、\textbf{注釈付きタグが推奨}されます。理由は以下のとおりです。

\begin{itemize}
  \item 作成者と日時が記録されるため、「誰がこのリリースを承認したか」が追跡できる
  \item メッセージに変更の要約を残せる
  \item \cmd{git describe} コマンドなどのツールが注釈付きタグを優先的に使用する
  \item 署名（GPG署名）を付与できるため、リリースの真正性を保証できる
\end{itemize}

軽量タグは、個人的な作業の目印や一時的なマーキングに使い、正式なリリースには注釈付きタグを使う、というのが一般的な使い分けです。

\subsection{タグの基本操作}

タグに関する基本的なGitコマンドを見ていきましょう。

\begin{lstlisting}[language=bash]
# 現在のHEADに注釈付きタグを作成
git tag -a v1.0.0 -m "Version 1.0.0 - Initial release"

# 過去のコミットにタグを付ける（コミットハッシュを指定）
git tag -a v0.9.0 -m "Beta release" abc1234

# タグの一覧を表示
git tag

# パターンで絞り込み
git tag -l "v1.*"

# タグの詳細情報を表示（注釈付きタグの場合、メッセージや作成者情報も表示）
git show v1.0.0

# タグをリモートにプッシュ（タグは通常のpushでは送信されない）
git push origin v1.0.0

# すべてのタグをリモートにプッシュ
git push origin --tags

# ローカルのタグを削除
git tag -d v1.0.0

# リモートのタグを削除
git push origin --delete v1.0.0
\end{lstlisting}

ここで重要なのは、\textbf{タグは \cmd{git push} だけではリモートに反映されない}という点です。通常のブランチのpushとは異なり、タグは明示的に \cmd{git push origin v1.0.0} または \cmd{git push origin --tags} で送信する必要があります。この仕様は、ローカルで一時的に付けたタグが意図せずリモートに送られることを防ぐための設計です。

\section{セマンティックバージョニング}

\subsection{MAJOR.MINOR.PATCH の意味}

ソフトウェアのバージョン番号は、自由に付けることもできますが、多くのプロジェクトでは\textbf{セマンティックバージョニング（Semantic Versioning、略称SemVer）}という標準的な規約に従っています。SemVerでは、バージョン番号を \cmd{MAJOR.MINOR.PATCH} の3つの数字で構成します。

\begin{lstlisting}[language={}]
v 2 . 3 . 1
  ↓   ↓   ↓
  |   |   └── PATCH: バグ修正（後方互換性あり）
  |   └────── MINOR: 新機能追加（後方互換性あり）
  └────────── MAJOR: 破壊的変更（後方互換性なし）
\end{lstlisting}

この3つの数字には明確な意味があり、バージョン番号を見るだけで「このアップデートは安全に適用できるか」を判断できるようになっています。

\subsection{各桁の判断基準}

\textbf{PATCH（パッチ）バージョン} --- バグ修正や小さな改善で、既存のAPIやインターフェースに変更がない場合に上げます。たとえば、計算結果の丸め誤差を修正した場合や、ドキュメントのtypoを修正した場合などです。ユーザーは安心してアップデートできます。

\begin{lstlisting}[language={}]
v1.2.3 → v1.2.4  （バグ修正のみ）
\end{lstlisting}

\textbf{MINOR（マイナー）バージョン} --- 新しい機能を追加したが、既存の機能との互換性は保たれている場合に上げます。新しいAPIエンドポイントの追加、新しいオプションパラメータの追加などが該当します。既存のコードを変更せずにアップデートできることが保証されます。MINORを上げた場合、PATCHは0にリセットします。

\begin{lstlisting}[language={}]
v1.2.4 → v1.3.0  （新機能追加、既存機能は維持）
\end{lstlisting}

\textbf{MAJOR（メジャー）バージョン} --- 既存のAPIやインターフェースに互換性のない変更（破壊的変更）を加えた場合に上げます。関数のシグネチャの変更、設定ファイルの形式変更、依存関係の大幅な変更などが該当します。ユーザーはアップデートの際にコードの修正が必要になる可能性があります。MAJORを上げた場合、MINORとPATCHは0にリセットします。

\begin{lstlisting}[language={}]
v1.3.0 → v2.0.0  （互換性のない変更あり）
\end{lstlisting}

\subsection{プレリリースバージョン}

正式リリースの前にテスト版を公開する場合は、バージョン番号にサフィックスを付けます。

\begin{tabular}{|l|l|l|}
\hline
\textbf{サフィックス} & \textbf{意味} & \textbf{安定度} \\
\hline
\cmd{v1.0.0-alpha} & アルファ版。初期テスト段階。機能が不完全な可能性あり & 低 \\
\hline
\cmd{v1.0.0-alpha.2} & アルファ版の2回目のリリース & 低 \\
\hline
\cmd{v1.0.0-beta} & ベータ版。機能は揃っているがバグが残っている可能性あり & 中 \\
\hline
\cmd{v1.0.0-beta.3} & ベータ版の3回目のリリース & 中 \\
\hline
\cmd{v1.0.0-rc.1} & リリース候補（Release Candidate）。重大なバグがなければ正式版になる & 高 \\
\hline
\end{tabular}

プレリリース版はSemVerの仕様上、同じ \cmd{MAJOR.MINOR.PATCH} の正式版よりも常に前（低い優先度）とされます。つまり、\cmd{v1.0.0-alpha < v1.0.0-beta < v1.0.0-rc.1 < v1.0.0} という順序です。

\subsection{v を付けるか付けないか}

タグに \cmd{v} プレフィックスを付けるかどうか（\cmd{v1.0.0} vs \cmd{1.0.0}）は、プロジェクトによって流儀が異なります。SemVerの仕様自体は \cmd{v} を含みませんが、GitHubでは \cmd{v} を付ける慣習が広く採用されています。GitHub CLIの \cmd{gh release create} も \cmd{v} 付きのタグを前提としている部分があります。大事なのはプロジェクト内で統一することです。

\section{リリースの作成}

GitHubの\textbf{リリース（Release）}は、タグの上に追加情報を載せたものです。タグが「コミットに付けた名前」なのに対し、リリースは以下の情報を含みます。

\begin{itemize}
  \item \textbf{リリースノート} --- このバージョンで何が変わったかの説明（変更履歴）
  \item \textbf{アセット（添付ファイル）} --- コンパイル済みバイナリ、インストーラー、ZIPファイルなどの配布物
  \item \textbf{メタ情報} --- ドラフトか公開済みか、プレリリースか正式版か
\end{itemize}

リリースを作ることで、ユーザーは「現在の安定版はどれか」「どのバージョンからどんな変更があったか」を簡単に確認できます。

\subsection{Web UIでの作成}

\begin{enumerate}
  \item リポジトリのメインページ右側にある「\textbf{Releases}」セクション（またはタグアイコン）をクリック
  \item 「\textbf{Create a new release}」ボタンをクリック（既存のリリースがある場合は「Draft a new release」）
  \item 「\textbf{Choose a tag}」で既存のタグを選択するか、新しいタグ名を入力して作成する
  \item \textbf{リリースタイトル}を入力（通常はバージョン番号、例: \cmd{v1.0.0}）
  \item \textbf{リリースノート}を記述する。「\textbf{Generate release notes}」ボタンをクリックすると、前回のリリースからの変更をPRベースで自動生成してくれる
  \item 必要に応じて、バイナリファイルやアーカイブを「Attach binaries」エリアにドラッグ\&ドロップ
  \item プレリリースの場合は「\textbf{Set as a pre-release}」にチェック
  \item 「\textbf{Publish release}」をクリック
\end{enumerate}

\screenshot{リリース作成画面。タグ選択、タイトル入力、リリースノート記入欄、Generate release notesボタンが見える状態}

「Generate release notes」機能は非常に便利です。前回のリリース以降にマージされたPRのタイトルと作成者を自動的にリストアップしてくれるため、手動で変更履歴を書く手間が大幅に削減されます。

\subsection{CLIでの作成}

GitHub CLIを使えば、コマンド1つでリリースを作成できます。

\begin{lstlisting}[language=bash]
# 基本的なリリース作成（タグがまだなければ自動作成される）
gh release create v1.0.0 --title "v1.0.0" --notes "Initial release"

# リリースノートを自動生成
gh release create v1.0.0 --generate-notes

# タイトルも自動（タグ名がそのままタイトルになる）
gh release create v1.0.0 --generate-notes --title "v1.0.0 - First Stable Release"

# ドラフトリリース（まだ公開しない）
gh release create v1.0.0 --draft --generate-notes

# プレリリース
gh release create v1.0.0-beta.1 --prerelease --generate-notes

# ファイルを添付してリリース
gh release create v1.0.0 \
  --generate-notes \
  ./dist/app-linux-amd64 \
  ./dist/app-darwin-arm64 \
  ./dist/app-windows-amd64.exe
\end{lstlisting}

リリースの管理に使えるその他のコマンドも見ておきましょう。

\begin{lstlisting}[language=bash]
# リリース一覧の表示
gh release list

# 特定のリリースの詳細表示
gh release view v1.0.0

# リリースのアセット（添付ファイル）をダウンロード
gh release download v1.0.0

# 特定のアセットだけダウンロード
gh release download v1.0.0 --pattern "*.tar.gz"

# リリースを削除
gh release delete v1.0.0

# ドラフトリリースを公開に変更
gh release edit v1.0.0 --draft=false
\end{lstlisting}

\subsection{ドラフトリリースの活用}

リリースを作成する際、すぐに公開せず\textbf{ドラフト}として保存しておくことができます。ドラフトリリースは、チームメンバーによるリリースノートのレビューや、QAテストの完了を待ってから公開する、といったワークフローで活用されます。

ドラフトリリースはリポジトリの Releases ページに表示されますが、「Draft」というバッジが付き、公開リリースの一覧には含まれません。準備が整ったら「Edit」から「Publish release」をクリック（またはCLIで \cmd{gh release edit v1.0.0 --draft=false}）するだけで公開できます。

\section{リリースのカスタマイズ}

\subsection{リリースノート自動生成のカテゴリ分け}

「Generate release notes」で自動生成されるリリースノートは、デフォルトではすべてのPRが1つのリストに並びます。プロジェクトが大きくなると、「新機能」「バグ修正」「ドキュメント」などのカテゴリに分けて表示したくなります。

これを実現するのが \cmd{.github/release.yml} ファイルです。このファイルでは、PRのラベルに基づいてリリースノートをカテゴリ分けするルールを定義します。

\begin{lstlisting}[language={}]
# .github/release.yml
changelog:
  exclude:
    labels:
      - "skip-changelog"       # このラベルが付いたPRはリリースノートから除外
    authors:
      - "dependabot"           # dependabotによるPRも除外
  categories:
    - title: "New Features"
      labels:
        - "enhancement"
        - "feature"
    - title: "Bug Fixes"
      labels:
        - "bug"
        - "bugfix"
    - title: "Documentation"
      labels:
        - "documentation"
        - "docs"
    - title: "Maintenance"
      labels:
        - "chore"
        - "dependencies"
        - "refactor"
    - title: "Other Changes"
      labels:
        - "*"                  # 上記のいずれにも該当しないPR
\end{lstlisting}

この設定を行うと、「Generate release notes」で生成されるリリースノートは以下のようにカテゴリ分けされます。

\begin{lstlisting}[language={}]
## What's Changed

### New Features
* Add dark mode support by @alice in #42
* Implement export to PDF by @bob in #45

### Bug Fixes
* Fix login redirect loop by @charlie in #43
* Fix date parsing in non-English locales by @alice in #47

### Documentation
* Update API reference by @dave in #44

### Other Changes
* Update CI workflow by @eve in #46

**Full Changelog**: https://github.com/owner/repo/compare/v1.0.0...v2.0.0
\end{lstlisting}

ポイントは、\textbf{ラベルの活用がリリースノートの品質に直結する}ということです。第11章で学んだラベルの運用が、ここで生きてきます。PRを作成する際にラベルを適切に付けていれば、リリース時にはボタン1つで整理されたリリースノートが手に入ります。

\subsection{exclude セクション}

\cmd{exclude} セクションでは、リリースノートに含めたくないPRを指定します。典型的なのは以下のケースです。

\begin{itemize}
  \item \textbf{dependabot}（依存関係の自動更新bot）によるPR。大量に発生するため、リリースノートが見づらくなる
  \item \textbf{skip-changelog}ラベルを付けた、リリースノートに記載する必要のない内部的な変更
\end{itemize}

\cmd{exclude} に指定されたPRは、どのカテゴリにもリストされなくなります。

\section{まとめ}

この章では、ソフトウェアをユーザーに届けるための「リリース管理」の仕組みを学びました。

\begin{itemize}
  \item \textbf{タグ}は特定のコミットに名前を付ける機能であり、リリースには\textbf{注釈付きタグ}を使うことが推奨されます
  \item \textbf{セマンティックバージョニング}（\cmd{MAJOR.MINOR.PATCH}）は、バージョン番号にルールを設けることで、アップデートの影響範囲をユーザーが判断できるようにする規約です
  \item \textbf{GitHubリリース}は、タグにリリースノートとアセットを追加したもので、Web UIとCLIの両方から作成できます
  \item 「\textbf{Generate release notes}」機能を使えば、PRベースで変更履歴を自動生成できます
  \item \cmd{.github/release.yml} を設定すると、\textbf{PRのラベルに基づいてリリースノートをカテゴリ分け}でき、ユーザーにとって見やすいリリースノートを作成できます
\end{itemize}

リリース管理は、開発チームとユーザーの間の\textbf{コミュニケーション手段}です。「何が変わったか」「安全にアップデートできるか」「過去のどのバージョンを使えばよいか」――これらの疑問に答えるのが、適切に管理されたタグとリリースノートの役割です。

また、第12章で学んだGitHub Actionsと組み合わせることで、タグをプッシュしたら自動でビルド・テスト・リリース作成が行われる、という完全自動化されたリリースパイプラインを構築することも可能です。
